\input eplain
\pdfpagewidth=148mm
\pdfpageheight=210mm
\hsize=128mm
\vsize=190mm

\pdfhorigin=0pt
\pdfvorigin=0pt
\hoffset=10mm
\voffset=10mm
\noindent grind \leaders\hbox to 10pt {\hss.\hss}  \hfill \xref{node5}
\smallskip
\noindent hw \leaders\hbox to 10pt {\hss.\hss}  \hfill \xref{node3}
\smallskip
\noindent news \leaders\hbox to 10pt {\hss.\hss}  \hfill \xref{node4}
\smallskip
\noindent notes \leaders\hbox to 10pt {\hss.\hss}  \hfill \xref{node0}
\smallskip
\noindent review \leaders\hbox to 10pt {\hss.\hss}  \hfill \xref{node1}
\smallskip
\noindent study \leaders\hbox to 10pt {\hss.\hss}  \hfill \xref{node2}
\smallskip
\vfill \break
{\medskip\noindent\bf\font\titlefont=cmss17 at 20pt\titlefont Timeline}\medskip
{\noindent \bf 2026-01-20:} notes\smallskip
{\noindent \bf 2026-01-22:} notes\smallskip
{\noindent \bf 2026-01-27:} review\smallskip
{\noindent \bf 2026-01-28:} hw, news, study\smallskip
{\noindent \bf 2026-01-29:} hw\smallskip
{\noindent \bf 2026-01-30:} grind\smallskip
{\noindent \bf 2026-02-01:} news\smallskip
{\noindent \bf 2026-02-03:} hw\smallskip
{\noindent \bf 2026-02-04:} hw\smallskip
{\noindent \bf 2026-02-06:} grind\smallskip
{\noindent \bf 2026-02-07:} news\smallskip
\vfill \break
\xrdef{node0}
{\medskip\noindent\bf\font\titlefont=cmss17 at 20pt\titlefont notes}\medskip
{\medskip\noindent\bf\font\titlefont=cmitt12 at 17pt\titlefont 2026-01-20}\medskip
{\smallskip\noindent\bf\font\titlefont=cmr12 at 13pt\titlefont $\bullet$ 16:54: MA123 note re-writing}\smallskip
Still trying to figure out a workflow for this class, but I found myself wanting to take my scribbles in class and rewrite them using my knowledge graph format.
\par

{\medskip\noindent\bf\font\titlefont=cmitt12 at 17pt\titlefont 2026-01-22}\medskip
{\smallskip\noindent\bf\font\titlefont=cmr12 at 13pt\titlefont $\bullet$ 16:50: MA 123 notes transcription}\smallskip
{\smallskip\noindent\bf\font\titlefont=cmr12 at 13pt\titlefont $\bullet$ 17:37: notes transcription}\smallskip
{\smallskip\noindent\bf\font\titlefont=cmr12 at 13pt\titlefont $\bullet$ 18:00: creating the graph}\smallskip
{\smallskip\noindent\bf\font\titlefont=cmr12 at 13pt\titlefont $\bullet$ 18:17: creating the graph}\smallskip
\xrdef{node1}
{\medskip\noindent\bf\font\titlefont=cmss17 at 20pt\titlefont review}\medskip
{\medskip\noindent\bf\font\titlefont=cmitt12 at 17pt\titlefont 2026-01-27}\medskip
{\smallskip\noindent\bf\font\titlefont=cmr12 at 13pt\titlefont $\bullet$ 16:35: review: complex conjugates, factoring (quickly)}\smallskip
One of the examples in the textbook uses a conjugate to find and equivalent limit that to apply limit laws with. I completely forgot what a conjugate was, but found a wiki page on the subject (specifically related to square roots).
\par

It turns out we are factoring, so I had to run home and remember some of the tricks involved in factoring quadratics. It's coming back.
\par

\xrdef{node2}
{\medskip\noindent\bf\font\titlefont=cmss17 at 20pt\titlefont study}\medskip
{\medskip\noindent\bf\font\titlefont=cmitt12 at 17pt\titlefont 2026-01-28}\medskip
{\smallskip\noindent\bf\font\titlefont=cmr12 at 13pt\titlefont $\bullet$ 10:20: some problem sets from chapter 2}\smallskip
Had some issues accessing the HW today, so I took some time this morning to do some problems from the book.
\par

I've also begun a "scratch" index system. Basically, I'm going to hold onto and index the scraps I work out problems on. Might be useful dunno yet.
\par

\xrdef{node3}
{\medskip\noindent\bf\font\titlefont=cmss17 at 20pt\titlefont hw}\medskip
{\medskip\noindent\bf\font\titlefont=cmitt12 at 17pt\titlefont 2026-01-28}\medskip
{\smallskip\noindent\bf\font\titlefont=cmr12 at 13pt\titlefont $\bullet$ 14:30: transcribe hw1 problems}\smallskip
{\smallskip\noindent\bf\font\titlefont=cmr12 at 13pt\titlefont $\bullet$ 15:00: more transcribing of hw1 problems}\smallskip
{\medskip\noindent\bf\font\titlefont=cmitt12 at 17pt\titlefont 2026-01-29}\medskip
{\smallskip\noindent\bf\font\titlefont=cmr12 at 13pt\titlefont $\bullet$ 09:03: problem sets hw1}\smallskip
2.1.1, 2.1.3, 2.1.4, 2.1.9, 2.1.10, 2.1.13
\par

{\smallskip\noindent\bf\font\titlefont=cmr12 at 13pt\titlefont $\bullet$ 11:58: hw1 psets}\smallskip
2.1.14
\par

{\smallskip\noindent\bf\font\titlefont=cmr12 at 13pt\titlefont $\bullet$ 17:10: hw1 psets}\smallskip
2.1.17, 2.2.2. The rest of the problemsets require looking at a supplied graph (I did not transcribe these).
\par

{\medskip\noindent\bf\font\titlefont=cmitt12 at 17pt\titlefont 2026-02-03}\medskip
{\smallskip\noindent\bf\font\titlefont=cmr12 at 13pt\titlefont $\bullet$ 16:24: writing out homework 2 problems by hand (part 1)}\smallskip
{\medskip\noindent\bf\font\titlefont=cmitt12 at 17pt\titlefont 2026-02-04}\medskip
{\smallskip\noindent\bf\font\titlefont=cmr12 at 13pt\titlefont $\bullet$ 12:22: homework 2 work}\smallskip
\xrdef{node4}
{\medskip\noindent\bf\font\titlefont=cmss17 at 20pt\titlefont news}\medskip
{\medskip\noindent\bf\font\titlefont=cmitt12 at 17pt\titlefont 2026-01-28}\medskip
{\smallskip\noindent\bf\font\titlefont=cmr12 at 13pt\titlefont $\bullet$ 18:59: logging, seting up some new things related to my calculus class}\smallskip
{\medskip\noindent\bf\font\titlefont=cmitt12 at 17pt\titlefont 2026-02-01}\medskip
{\smallskip\noindent\bf\font\titlefont=cmr12 at 13pt\titlefont $\bullet$ 10:22: scanning first set of note cards as an experiment}\smallskip
Not sure where to put these, I just made a folder on my computer
\par

{\smallskip\noindent\bf\font\titlefont=cmr12 at 13pt\titlefont $\bullet$ 10:35: scanning scratch sheets as an experiment}\smallskip
Putting them in the same place as the cards.
\par

{\smallskip\noindent\bf\font\titlefont=cmr12 at 13pt\titlefont $\bullet$ 11:12: inventing some kind of ad-hoc annotation system literally right now}\smallskip
It's clear to me that if I don't put things into bins and add annotations, it's going to get lost.
\par

I'm just going to open up ed and start typing notes.
\par

{\medskip\noindent\bf\font\titlefont=cmitt12 at 17pt\titlefont 2026-02-07}\medskip
{\smallskip\noindent\bf\font\titlefont=cmr12 at 13pt\titlefont $\bullet$ 10:57: I failed yesterdays quiz}\smallskip
After thinking about it, I think it's because I took too long. The problem was immediately familiar, based on a homework problem (based on a textbook problem), but I just couldn't do the alegebra fast enough (10 minutes for two problems).
\par

So, I'm coming up with a new strategy. Don't write up the homework problems by hand. Type them out in TeX, and in a way that they can be easily parsed and turned into divisible chunks of data later. On friday, practice the homework problems instead of textbook problems. Make sure there are answers by then. Work in a routine of reviewing problems from previous weeks somehow.
\par

\xrdef{node5}
{\medskip\noindent\bf\font\titlefont=cmss17 at 20pt\titlefont grind}\medskip
{\medskip\noindent\bf\font\titlefont=cmitt12 at 17pt\titlefont 2026-01-30}\medskip
{\smallskip\noindent\bf\font\titlefont=cmr12 at 13pt\titlefont $\bullet$ 11:42: textbook pset grinding}\smallskip
2.2.19, 2.2.7, 2.2.4, 2.2.5, 2.2.6, 2.2.18
\par

{\medskip\noindent\bf\font\titlefont=cmitt12 at 17pt\titlefont 2026-02-06}\medskip
{\smallskip\noindent\bf\font\titlefont=cmr12 at 13pt\titlefont $\bullet$ 11:48: grinding}\smallskip
2.3 ex 2, ex 4, ex 7, 2.3.29, 2.3.37, 2.3.55, 2.3.67, 2.3.72 2.4 definition of infinite limits
\par

\bye
